% !TEX root = ../../../main.tex
% Appendix entry: interrogation.intervals -- "Intervals" (Appendix A).
\entryfield{Algorithm}
The adapter (\code{_run} and \code{measure} in \code{Adapters/interrogation_intervals.py})
calls the one interval implementation in the repository,
\code{Working/interrogation/intervals.py} (\code{inter_event_intervals},
\code{interval_statistics}). It never reads the signal.

\begin{enumerate}
\item \emph{The anchor.} With \emph{anchor} = \texttt{onset} and an upstream
  \code{onset_idx} column, each event's anchor $a_j$ is its onset sample (from
  \blockref{interrogation.event_shape}: the foot of the fall, walked back to the shoulder it
  departs from). Otherwise $a_j$ is the span start. The choice is made for the whole input,
  not per event: with the column present, an event whose onset is NaN is left out of the
  timing; it does not fall back to its span start. \code{meta.anchor_used} says which.
\item \emph{The groups.} The key of event $j$ is its upstream \code{group} value as text, or
  \texttt{all} when there is no \code{group} column.
\item \emph{Per group}, in order of first appearance: keep the events with a finite anchor,
  sort them by anchor (stable), convert to seconds $t_k = a_k / f_s$, and form the gaps
  $I_k = t_{k+1} - t_k$, $k = 1,\dots,m-1$ for $m$ events.
  \code{interval_before_s} of the $(k{+}1)$th event is $I_k$ and \code{interval_after_s} of
  the $k$th is $I_k$; the group's first event has no gap before and its last none after
  (NaN).
\item \emph{Statistics} of the group's finite gaps (\code{interval_statistics}), with
  $\bar I$ their mean and $n = m-1$ their number:
  \begin{itemize}
  \item \code{median_s}, \code{mean_s}, \code{min_s}, \code{max_s};
  \item $\mathrm{CV} = s/\bar I$ with the sample SD $s$ (ddof 1) (\code{regularity},
    moved from \code{sequences11});
  \item \code{r2_trend}: $r^2$, the squared Pearson correlation of $I_k$ with its index $k$;
    0 for a perfectly constant gap, where $r$ is undefined;
  \item \code{drift_ratio}: the mean of the last $\lfloor n/3 \rfloor$ gaps over the mean of
    the first $\lfloor n/3 \rfloor$ (at least one gap each); 1 is no drift, above 1 the gaps
    widen;
  \item the local measures over neighbouring pairs, $r_k = |I_k - I_{k+1}|/(I_k + I_{k+1})$,
    leaving out pairs whose two gaps are both 0\,s (a 0/0):
    \[ \mathrm{LV} = \frac{3}{p}\sum_{k} r_k^2 , \qquad
       \mathrm{CV2} = \frac{2}{p}\sum_{k} r_k , \]
    with $p$ the number of kept pairs (\code{n_local_pairs}). LV is
    \citet{shinomoto2003differences}'s local variation (the 2009 form, eq.~2, which the
    repository's sources verified); CV2 is \citet{holt1996comparison}'s. Both are about 0
    for clockwork and 1 for Poisson timing. The refractory-corrected LvR is not offered: its
    constant (5\,ms) was fitted to cortical neurons.
  \end{itemize}
\end{enumerate}

\code{inter_event_intervals} sorts before differencing, so an unsorted input never gives a
negative gap. The block applies no threshold to these numbers: it reports them, and the
judgement of whether a run is a regular train is left to the reader.

\entryfield{Assumptions}
\begin{itemize}
\item \emph{Units.} Gaps are in seconds, from sample indices divided by the chain's $f_s$.
  The signal's unit never enters. CV, $r^2$, the drift ratio, LV and CV2 have no unit.
\item \emph{Index frame.} Starts and \code{onset_idx} are both span-relative (counted from
  the start of the span the chain ran on), so their differences are the same in any frame.
\item \emph{What is timed.} Onset to onset (or span start to span start). The block does
  not measure the quiet time between one event's end and the next's start, nor peak to
  peak.
\item \emph{Missing values.} The trace is not read, so a NaN stretch is invisible: the gap
  across it is reported as one long gap. A NaN anchor drops the event from its group's
  timing (step 1).
\item \emph{Edges.} An event near either end of the span is treated like any other; the gap
  to an event outside the span is unknown, so the first and last events of each group have
  one NaN gap each.
\item \emph{How many events.} One event: no gap, every statistic \code{None}. Two events:
  one gap, median and mean but no CV, no $r^2$, no drift ratio and no LV or CV2 (all
  \code{None}). LV and CV2 need three events. Every undefined value is \code{None}, never 0, inf or
  NaN, and \code{local_note} says why.
\end{itemize}

\entryfield{Outputs}
A \texttt{SpanSet} with the input's spans, labels, scores and feature columns unchanged, plus
\code{interval_before_s} and \code{interval_after_s} (seconds; replacing columns of the same
name upstream). \code{meta.interval_stats} maps each group key to \code{n_events},
\code{n_intervals}, \code{median_s}, \code{mean_s}, \code{min_s}, \code{max_s}, \code{cv},
\code{r2_trend}, \code{drift_ratio}, \code{lv}, \code{cv2}, \code{n_local_pairs} and
\code{local_note}. \code{meta} also holds \code{rules} (the printed rule for every
quantity), \code{anchor_used}, \code{n_events}, \code{n_groups} and
\code{n_groups_local_unmeasured} (groups with no LV).

A group key written by an upstream feature block can carry a suffix:
\texttt{\textasciitilde noise\_like} or \texttt{\textasciitilde untested} (from
\blockref{interrogation.event_families} and \blockref{interrogation.cycle_shapes}). Such a
row is the rhythm of a family that did not pass, or was not given, a noise test; the key
\texttt{-1} collects events in no train or family, and its gaps are between unrelated
events.

\entryfield{Cost}
$O(n \log n)$ in the number of events, dominated by the per-group sort; negligible memory.
The span length enters only through the number of events. Not meant for the cluster.

\entryfield{Failure modes}
\begin{itemize}
\item \emph{Duplicate onsets.} One duplicated onset in an otherwise regular train (gaps 30,
  30, 0, 30, 30, 30) gives LV 1.20 and CV2 0.80, which reads as random (module docstring).
\item \emph{Few gaps.} LV's Poisson SD is 0.89 at 2 gaps, 0.59 at 4 and 0.37 at 9 (module
  docstring); \citet{shinomoto2003differences} used fragments of at least 100 gaps.
\item \emph{The sampling grid.} Onsets fall on whole samples. A strictly periodic 2.5\,s
  train on a 1\,Hz grid reads CV2 0.20; a small LV or CV2 for gaps of a few samples is grid
  noise.
\item \emph{CV and drift.} CV mixes a slow change of rate with irregularity: in the tests a
  ten-fold rise in rate gives CV above 1.15 while LV and CV2 stay near 1. Read CV beside
  \code{drift_ratio} and \code{r2_trend}.
\item \emph{A missed event doubles a gap; a split event halves one.} The block trusts its
  input; a detector's misses show up here as outliers, not as warnings.
\item \emph{Pooling by mistake.} Without a \code{group} column every event in the span is
  one group. That is right for one channel's span, and wrong if events of different kinds
  share the span and you meant to separate them: put Trains or Event families with
  \emph{write\_group} on before it.
\end{itemize}

\entryfield{Validation}
The tests check numbers, not only shapes. \code{tests/test_adapter_intervals.py}: gaps
worked by hand from span starts (10, 15, 5\,s at 10\,Hz) and from an upstream onset column;
CV, $r^2$ and drift equal to the hand-computed values; groups never pooled; unsorted spans
ordered first; one event giving no statistics and no error; a regular 3\,s train giving CV 0.
\code{tests/test_intervals_lv.py}: LV and CV2 equal to the published formulas on a small
case; 0 for a regular train; within 0.1 of 1 for 5{,}000 Poisson gaps; CV fooled by a
ten-fold rate rise while LV and CV2 are not; the handling of zero-gap pairs; per-group
reporting differing from pooled. No test compares the block's output with intervals measured
by hand on a fungal recording.

\entryfield{References}
\citet{shinomoto2003differences} (LV); \citet{holt1996comparison} (CV2). CV, $r^2$ and the
drift ratio are the project's own regularity gate from \code{sequences11}, now in
\code{Working/interrogation/intervals.py}. Background:
\code{docs/research/analysis-blocks-primary-sources.md}, section 9, and
\code{docs/research/analysis-blocks-from-literature.md}.

\entryfield{Tests}
\code{tests/test_adapter_intervals.py} (the block, and that the three earlier interval
implementations now share \code{inter_event_intervals}); \code{tests/test_intervals_lv.py}
(LV and CV2); \code{tests/test_units_reach_the_blocks.py},
\code{tests/test_template_event_trains.py}, \code{tests/test_template_event_families.py},
\code{tests/test_template_chain_cycles.py} and \code{tests/test_end_to_end.py} run it inside
chains, with per-group rhythms checked in the families template (within 10\,\%) and the cycle
shapes template.
